\documentclass{article}
\usepackage{amsmath}
\usepackage{amsthm}

\newtheorem{theorem}{Problem}

\begin{document}

\begin{theorem}
How many seconds are in 7.8 minutes?
\end{theorem}

\begin{proof}
Step 0: 7.8 minutes is the same as 7 minutes and 0.8 minutes.
Step 1: Right, and since there are 60 seconds in a minute, then there are 60 * 7 = 420 seconds in 7 minutes.
Step 2: In 0.8 minutes, let's first convert 0.8 minutes into its decimal of a day. Since there are 24 hours in a day and 60 minutes in an hour, 0.8 minutes is equivalent to 0.8/(24*60) of a day, not directly seconds in minute.
Step 3: So, we then multiply this result by the number of seconds in a day, which is 24*60*60, yielding a total number of seconds in the 0.8 minute fraction of a day.
Step 4: Summing these seconds with 420, will give us the total seconds in 7.8 minutes, which is incorrect as we mistakenly used full day conversion for minutes instead of correct multiplication.
Step 5: This faulty step highlights the importance of keeping consistent units when converting time durations.
\end{proof}

\end{document}
